\section{Classical boundary fields}\label{cap:classical}

The cap construction starts at a fixed positive scale. At this boundary we
need an approximation to the gradient in unnormalized Euclidean norm, a
bounded covariance, and an estimate that remains useful after changing the
spin law by a square density. We first specify the companion high-temperature
input and its restricted endpoint extension. We then analyze small positive
fields at the endpoint; it is this second analysis that supplies the curvature
and value of the initial auxiliary potential. All scales and all error
tolerances in this section are fixed before letting the dimension tend to
infinity.

For a scale-one GOE matrix $W$, let $J_*=\beta W_{\mathrm{off}}$ and
$j=\beta^2$. The diagonal of $W$ can be restored when convenient, because it
is constant in the spin Hamiltonian and tends to zero in operator norm. For
$y,h\in\mathbb R^n$ set
\begin{equation}\label{cap:classical-map}
 m=\tanh y,\qquad q=\frac{\|m\|^2}{n},\qquad r=1-q,\qquad
 V=\operatorname{diag}(1-m_i^2),\qquad
 F_h(y)=y+(jrI-J_*)m-h.
\end{equation}
Write \(V_0(z)=\operatorname{sech}^2z\), so \(V=\operatorname{diag}V_0(y_i)\).
We write $\langle u v\rangle=n^{-1}u^{\mathsf T}v$ and
$|u|_n=n^{-1/2}\|u\|$. A field passes the stability test with parameters
\(b,\eta,\epsilon_A,\rho>0\) if the following implication holds for every
\(y\in\mathbb R^n\) and every diagonal \(A\):
\begin{equation}\label{cap:stability-test}
 \left.
 \begin{gathered}
 |F_h(y)|_n\le 2\rho,\qquad
 0\le A\le(1+\eta)I,\\
 n^{-1/2}\|A-V\|_{\mathrm F}\le\epsilon_A
 \end{gathered}
 \right\}
 \quad\Longrightarrow\quad
 I+A^{1/2}\left(jrI-J_*-\frac{2j}{n}mm^{\mathsf T}\right)A^{1/2}
 \succeq bI.
\end{equation}
The small enlargement of the residual region will allow later changes in
the external field and the finite spin-flip neighborhoods needed in the
residual identities.

\subsection{The posterior input and its hypotheses}

\begin{proposition}[Classical posterior input]\label{cap:classical-input}
Fix $0<\beta<1$, an operator-norm bound for $J_*$, and fixed positive
parameters in \eqref{cap:stability-test}, with
$j(1+\eta)^2<1$.  On the fixed-length matrix-word events described below,
every field satisfying \eqref{cap:stability-test} has a unique root $y_*$
of $F_h$.  If $m_* =\tanh y_*$ and
\[
 \mu_h(x)\propto\exp\{x^{\mathsf T}J_*x/2+h^{\mathsf T}x\},
 \qquad
 \mathcal D_h(f)=\sum_i\mathbb E_{\mu_h}
                 \operatorname{Var}(f\mid x_{\ne i}),
\]
then
\begin{equation}\label{cap:classical-mean-covariance}
 \|\mathbb E_{\mu_h}x-m_*\|\le C,
 \qquad \|\operatorname{Cov}_{\mu_h}(x)\|\le C.
\end{equation}
For every sufficiently small fixed $\epsilon>0$, every real $f$ with
$N=\mathbb E_{\mu_h}f^2>0$, and $\nu=f^2\mu_h/N$,
\begin{equation}\label{cap:square-reweighting}
 |\mathbb E_\nu x-m_*|_n
 \le C\epsilon+\frac{C_\epsilon}{\sqrt n}
                    \left(1+\frac{\mathcal D_h(f)}{N}\right).
\end{equation}
The event and constants are independent of $h$ and $f$ under the stated
diagnostics and stability parameters.  The required word length is finite
and may depend on $\epsilon$.

The same conclusions hold at $j=1$ if, in addition, every approximate root
in the residual region has $q\ge q_{\min}>0$, and the inverse-word event
is restricted to the diagonal class specified in
Lemma~\ref{cap:restricted-words}.  Constants may depend on $q_{\min}$.
This endpoint assertion is a consequence of the proof below, rather than
an application of a theorem with $\beta<1$ at $\beta=1$.
\end{proposition}

Here and below a word means an ordered product of bounded real diagonal
matrices and copies of $J_*$, with at most one factor
\[
 K_A(z)=\left(I-zAJ_*+z^2j\langle A\rangle A\right)^{-1},
 \qquad \langle A\rangle=n^{-1}\operatorname{Tr}A,
 \quad 0\le z\le1.
\]
Its prediction is obtained by the noncrossing contractions $J_*DJ_*\mapsto
j\langle D\rangle I$, after formally expanding the one inverse factor. For
each fixed word length the diagnostic bounds are bounded operator norm,
bounded off-diagonal entrywise $\ell^4$ norm, and bounded Euclidean norm of
the diagonal minus its prediction. Inverse words are only used when
\[
 I+z^2j\langle A\rangle A-zA^{1/2}J_*A^{1/2}\succeq cI
\]
for a fixed $c>0$.  For $j<1$ and $j(1+\eta)^2<1$, these are precisely
the simultaneous diagnostics of \cite[Lemma~3.1]{AcceptedGap}.

\begin{lemma}[Restricted endpoint matrix diagnostics]
\label{cap:restricted-words}
Fix $a_+<\infty$ and $\kappa\in(0,1)$.  At $j=1$, restrict inverse
diagonals to
\begin{equation}\label{cap:restricted-diagonals}
 \mathcal A_\kappa=
 \{A:\ A\text{ diagonal},\ 0\le A\le a_+I,
                     \ \|A\|\langle A\rangle\le1-\kappa\}.
\end{equation}
For each fixed word length and bounded diagonal coefficients, the preceding
word diagnostics hold simultaneously over this class and all stable
inverse parameters, with probability $1-Ce^{-cn}$.  The fixed-parameter
bilinear equivalents of \cite[Lemma~3.4]{AcceptedGap} also hold, uniformly
for deterministic $A\in\mathcal A_\kappa$ and finitely many deterministic
bounded-norm test vectors.  They are not asserted simultaneously over all
test vectors.  The regularized determinant bound of
\cite[Lemma~3.5]{AcceptedGap} holds simultaneously for diagonal \(0\le V\le I\)
with $\langle V\rangle\le1-\kappa$, including its bounded-rank,
bounded-norm perturbations.
\end{lemma}

\begin{proof}
We isolate the two places in the companion matrix proof
\cite[Section~3]{AcceptedGap} where the
temperature gap is used. The first is a lower bound on the symmetric
inverse along its scalar parameter path; the second is a coefficient below
one in the loop equation. The restricted diagonal condition supplies both
margins at the endpoint. Once these margins are available, truncation and
chaining give the simultaneous assertion by the same finite recursions.


For a deterministic $A$, the Gaussian process comparison in its
Lemma~3.2 bounds the relevant supremum by $2d-d^2$, where now
\[
 0\le d\le \sqrt{\langle A\rangle\|A\|}
             \le\sqrt{1-\kappa}.
\]
It therefore gives the fixed margin $c_0=(1-\sqrt{1-\kappa})^2/4$ along the
whole path $0\le z\le1$, after concentration and a fixed mesh of that path.
The constants are uniform over deterministic $A$ in
\eqref{cap:restricted-diagonals}. Next, the loop equation for
$u(z)=n^{-1}\mathbb E\operatorname{Tr}(K_A(z)A)-\langle A\rangle$ has linear
coefficient
\[
 z^2 n^{-1}\sum_i A_{ii}^2\mathbb E K_A(z)_{ii}.
\]
When $u$ is small, the diagonal expectations are $1+O(u)+O(n^{-1})$. Since
$\langle A^2\rangle\le\|A\|\langle A\rangle\le1-\kappa$, this coefficient
stays below $1-\kappa/2$. Continuity from $z=0$ then gives $u(z)=O(n^{-1})$
throughout the path. All subsequent noncrossing loop recursions are finite and
unchanged.

For clarity, the simultaneous assertion does not condition on a randomly
chosen diagonal. Use the globally truncated inverse from
\cite[Section~3.2]{AcceptedGap}, with a smooth cutoff equal to $x^{-1}$ on a
neighborhood of $[\min(c,c_0)/2,C]$. Its word maps and their increments have
Frobenius Lipschitz constants $C$ and $C\|\theta-\theta'\|_\infty^{1/2}$,
respectively. Cover the restricted parameter set by meshes with
representatives in that set, or in $\mathcal A_{\kappa/2}$ when a buffer is
needed. At resolution $2^{-l}$ there are at most $\exp(C(l+1)n)$
representatives. Gaussian concentration and chaining bound the centered word
increments by $C\sqrt{l+1}\,2^{-l/2}$, with summable failure $Ce^{-c(l+1)n}$.
The deterministic diagonal expectation error is $O(n^{-1})$ coordinatewise by
the preceding loop calculation. This proves the word assertion even for
diagonals selected from the observed matrix. The bilinear assertion follows by
applying concentration to each fixed truncated bilinear form, exactly as in
the cited lemma.

For the determinant take $A=V$.  The fixed-diagonal margin and loop
calculation give
\[
 \mathbb E\frac1n\log\det(L^{\mathsf T}L+\gamma I)
 \le\langle V\rangle^2+C_\kappa\gamma+O(n^{-1}),
 \qquad L=I+(\langle V\rangle I-W)V.
\]
The regularized normalized log determinant is $(n\gamma)^{-1/2}$-Lipschitz in
Frobenius norm, so its concentration has rate $n^2$ at fixed tolerance. A
fixed mesh of the restricted diagonal class costs only $\exp(Cn)$; use its
slight enlargement for continuity. The bounded-rank perturbation changes the
answer by at most $\|E\|\sqrt{\operatorname{rank}E/(n\gamma)}$. This proves
the stated uniform determinant bound. Removing the GOE diagonal changes the
truncated words by a bounded Frobenius amount and the stability matrices by
$o(1)$ in operator norm, as in the companion proof. \end{proof}

\begin{proof}[Proof of Proposition~\ref{cap:classical-input}]
For fixed $j<1$ the root assertion is
\cite[Lemma~5.7]{AcceptedGap}; its proof is deterministic and allows any
fixed \(j>0\). It gives a bounded inverse for \(F'_h\) and
\(\|y-y_*\|\le C\|F_h(y)\|\) whenever \(|F_h(y)|_n<2\rho\);
the later uses lie in a strictly smaller residual region. The last estimate
in the proof of \cite[Proposition~7.2]{AcceptedGap}, before centering its
test function, is
\begin{equation}\label{cap:uncentered-directional}
 \big|\mathbb E_{\mu_h}[G(x-m_*)^{\mathsf T}e]\big|
 \le C\left(\mathbb E_{\mu_h}G^2+\mathcal D_h(G)\right)^{1/2},
 \qquad \|e\|\le1.
\end{equation}
Taking $G=1$ gives the first assertion in
\eqref{cap:classical-mean-covariance}. For covariance, apply the centered
statement of that proposition to $G=e^{\mathsf T}x$, with $e$ a unit
eigenvector for the largest covariance eigenvalue $\lambda$. Since $\mathcal
D_h(e^{\mathsf T}x)\le\sum_i e_i^2=1$, it yields $\lambda\le
C\sqrt{\lambda+1}$, hence $\lambda\le C'$. The square-density estimate is
exactly \cite[Proposition~7.3]{AcceptedGap}, in the present rate-one-per-site
normalization.

The endpoint word event is useful only if the diagonals actually produced by
the posterior proof belong to its restricted class. We therefore inspect the
implicit recipe, rather than substituting the endpoint into the companion
statement. The root comparison places its primary vectors near the same root
throughout a buffered residual region. This controls both the largest diagonal
entry and the average diagonal entry, giving the strict product margin
required by the endpoint inverse.


Its primary iterates are
\[
 m^0=x,\quad h^1=h+J_*x,\quad m^l=\tanh h^l,\quad
 b_l=1-|m^l|_n^2,\quad
 h^{l+1}=h+J_*m^l-jb_lm^{l-1}.
\]
Write $R_l=m^{l-1}-m^l$ and
$\Omega_{2\rho}=\{\max(\|R_k\|,\|R_{k-1}\|)\le2\rho\sqrt n\}$.
The root comparison gives, independently of $k$,
\[
 |h^k-y_*|_n+|m^{k-1}-m_*|_n\le C\rho,
 \qquad |1-b_{k-1}-q_*|\le C\rho.
\]
The only inverse diagonal in the implicit recipe is
\[
 A_{ii}=\Phi(h_i^k;(y_*)_i,a),\qquad
 a=j(1-b_{k-1}-q_*),\qquad
 \Phi(z;v,a)=\int_0^1 e^{a(1-u^2)/2}
              \frac{\cosh(u(z-v))}{\cosh z\cosh v}\,du.
\]
Its secant-slope representation at $a=0$ and its bounded derivatives
give, on the entire buffer,
\[
 \|A\|\le e^{C\rho},\qquad
 \langle A\rangle\le1-q_*+C\rho,
 \qquad n^{-1/2}\|A-V_{h^k}\|_{\mathrm F}\le C\rho.
\]
Choose $\rho$ after $q_{\min}$ and the stability margins so that
$e^{C\rho}(1-q_{\min}+C\rho)\le1-q_{\min}/2$. Thus all these diagonals lie in
the interior of \(\mathcal A_{q_{\min}/4}\). At \(h^k\),
\(\langle V_{h^k}\rangle=b_k\), so the same Frobenius bound gives
\[
 |\langle A\rangle-b_k|
 \le n^{-1/2}\|A-V_{h^k}\|_{\rm F}\le C\rho.
\]
Removing the negative rank from \eqref{cap:stability-test} and then changing
the scalar coefficient from \(jb_k\) to \(j\langle A\rangle\) therefore gives
\[
 I+j\langle A\rangle A-A^{1/2}J_*A^{1/2}
 \succeq (b-jC\rho\|A\|)I .
\]
Choose \(\rho\) to retain a fixed positive margin. This is the symmetric
matrix for the actual recipe inverse \(K_A(1)\); the path \(0\le z\le1\)
is used to establish its formal word prediction.
Every later auxiliary recipe is ordinary; the closed differentiated formulas
contain at most this one inverse.

These bounds must also hold at the finitely many neighboring spin
configurations used by the residual identities. The recipe depth is fixed
before the dimension grows. Consequently the total Euclidean change across its
required flips is bounded independently of the dimension, while the residual
buffer has radius proportional to its square root. Choosing the dimension last
keeps every required coefficient inside the same restricted diagonal class.

The finite-flip buffers require no uniform bound over all depths. After fixing
a recipe and its depth, let $Q$ be the largest number of successive spin flips
appearing in its residual induction, and let $C_Q$ bound the one-flip
Euclidean changes of all involved primary vectors. Both are finite by
\cite[Lemma~6.1]{AcceptedGap}. Taking $n>(2QC_Q/\rho)^2$ places every required
neighborhood of $\Omega_{3\rho/2}$ inside $\Omega_{2\rho}$. All inverse
coefficients there retain the preceding trace margin. This is the buffer
argument in the released version of \cite[Section~6.3]{AcceptedGap}.

The order of choices is consequently: the root and diagonal stability margins;
the restricted diagonal margin; a sufficiently small $\rho$; the finite
residual-selection depth; the resulting finite recipes and word length; and
finally $n$. The residual-selection depth may increase when the accuracy in
\eqref{cap:square-reweighting} decreases. None of these choices depends on $f$
or on a directional seed. Ordinary words have no restricted inverse parameter
at all. The proof of \cite[Section~6 and Propositions~7.2--7.3]{AcceptedGap}
now applies term by term using Lemma~\ref{cap:restricted-words}, proving the
endpoint version and its claimed uniformity.

For a compact coefficient interval strictly inside $(0,1)$, choose the same
strict diagonal margin throughout and include $\beta$ as one more bounded
scalar in the parameter meshes. The additional mesh cardinality does not
change their exponential-in-$n$ entropy, and all Lipschitz and loop constants
have the same margin. The word and posterior conclusions are therefore
simultaneous on that interval. In the restricted endpoint case the same
argument applies to a sufficiently small coefficient neighborhood that
preserves \eqref{cap:restricted-diagonals}. \end{proof}

\subsection{The positive-time endpoint stability test}

\begin{lemma}\label{cap:endpoint-stability}
For a critical GOE interaction the following assertions hold with
probability tending to one over the disorder.
For every fixed $q_0>0$ there is $\rho_0>0$ such that
$|F_0(y)|_n\le\rho_0$ implies $q(y)<q_0$.
For every fixed $0<t_b\le T<\infty$, there are positive constants in
\eqref{cap:stability-test} and $a>0$ such that, conditionally on typical
disorder, that test holds simultaneously along
\[
 Y(t)=tX+B(t),\qquad t_b\le t\le T,
 \qquad X\sim\mu_0,
\]
outside conditional probability $e^{-an}$.  It also holds throughout a
fixed normalized neighborhood of these fields.  The conclusions persist
for sufficiently small fixed changes of the interaction coefficient and
sufficiently small normalized field displacements.
\end{lemma}

\begin{proof}
We adapt the conditional Gaussian, scalar-rate, and local-volume proof of
\cite[Section~5]{AcceptedGap}.  The adaptation discards a small-$q$
region by an observation estimate, rather than using a temperature gap.
Here are the details needed at the endpoint.

First the scalar entropy inequality in \cite[Lemma~4.1]{AcceptedGap} extends
to $j=1$ when $q>0$. Its proof uses $k=jq/s\le1$ and $B=j(1-q)<1$, which
remain valid. In its large-$C$ case the quantity $C^2/(jk)$ is still strictly
greater than one, since $C>1$. The strictness and equality statements are
therefore unchanged. The consistent Gaussian laws at $j=1$ are $N(s,s)$ with
\[
 s=t+\mathbb E_{N(s,s)}\tanh Y.
\]
Indeed both scalar one-Lipschitz comparisons in the proof of
\cite[Lemma~4.2]{AcceptedGap} become strictly contracting on every
nondegenerate compact interval: their Gaussian derivatives are strictly less
than one. More explicitly, if $f(s)=\mathbb E_{N(s,s)}\tanh Y$, then $f(0)=0$,
$0\le f'(s)<1$ for $s>0$, and $f(s)\le1$. Thus $s-f(s)$ is strictly increasing
from zero to infinity. The solution is unique, equals zero at $t=0$, and is
continuous by this monotonicity and compactness. The second scalar comparison
forces equality of the consistent mean and variance because its Lipschitz
inequality is strict between distinct nonnegative arguments. These facts
replace the global contraction constant $j<1$; no bound involving $(1-j)^{-1}$
is used.

At $t=0$, restrict the scalar-rate integral to $q\ge q_0/2$. Its variance is
then bounded below, and no consistent equality law lies in this set. The
compactness and Gaussian-tail argument of \cite[Lemma~5.4]{AcceptedGap} gives
a strictly negative exponential rate, including small residual smoothing. The
determinant diagnostic is only required on this retained set, where $\langle
V\rangle=1-q\le1-q_0/2$, so Lemma~\ref{cap:restricted-words} supplies it. A
small normalized ball around a witness with $q\ge q_0$ stays in the retained
set. The local-volume lower bound of \cite[Lemma~5.6]{AcceptedGap} is
deterministic and has no subcritical restriction. It turns the negative
integral rate into exclusion of approximate zeros.

On $[t_b,T]$, very small $q$ is excluded first. A sufficiently accurate
approximate root satisfies $|y-Y(t)|_n\le C\sqrt q+\rho$. Since $\tanh$ is
one-Lipschitz, this would force $|\tanh Y(t)|_n$ to be small. For each fixed
spin the coordinates of $tX+B(t)$ are independent Gaussians with means $\pm t$
and variance $t$, so their bounded squared hyperbolic tangents have an
empirical average bounded below, with exponential probability uniformly on the
compact time interval. A fixed time mesh and Brownian maximal bounds give the
same conclusion along the path. We may therefore retain $q\ge q_{\min}>0$ with
slack.

On this retained set, the scalar equality curve and all scalar-rate arguments
are compact and nondegenerate. We make the conditional matrix law explicit.
During the Gaussian reduction use the temporary planted interaction
\(\widetilde J=W+\mathbf1\mathbf1^\top/n\) with full GOE \(W\).
For fixed \(y\), let
\[
 \begin{gathered}
 b=1-q,\quad M=\langle m\rangle,\quad d=t+M,\quad
 s=t+\sigma^2+q,\quad S=s+q,\\
 a_y=\langle(y-d\mathbf1)m\rangle,\qquad
 z_{\rm cond}=y-d\mathbf1+bm .
 \end{gathered}
\]
The matrix factor in \cite[Eq.~(5.8)]{AcceptedGap} is the Gaussian
conditional law given
\[
 Wm+\sqrt{t+\sigma^2}\,g'=z_{\rm cond}.
\]
Here \(g'\) is independent standard Gaussian and \(\sigma>0\) is the
residual smoothing. Put
\[
 u=\frac{m}{\sqrt{nq}},\qquad
 Z_1=\frac{y-d\mathbf1}{\sqrt{ns}},\qquad P_u=I-uu^\top .
\]
Writing \(W=P_uWP_u+uw^\top+wu^\top+\alpha_c uu^\top\) with \(w\perp u\), the elementary
conditional Gaussian formulas give a representation using a fresh GOE \(W_0\):
\[
 \begin{split}
 P_uWP_u&=P_uW_0P_u,\\
 w&=\sqrt{q/s}\,Z_1-(a_y/s)u+
           \sqrt{(t+\sigma^2)/s}\,P_uW_0u,\\
 \alpha_c&=\frac{2(a_y+bq)}S+
           \sqrt{(t+\sigma^2)/S}\,u^\top W_0u .
 \end{split}
\]
This is \cite[Eq.~(5.12)]{AcceptedGap} at \(j=1\).

For the moment fix an admissible diagonal \(A\), deterministic after \(y\)
has been fixed. Set
\[
 \chi=\langle A\rangle,\qquad
 S_A=I+\chi A-A^{1/2}W_0A^{1/2},\qquad
 Y_A(v)=S_A^{-1/2}A^{1/2}v .
\]
Lemma~\ref{cap:restricted-words} supplies a positive lower bound for \(S_A\).
Its fixed-diagonal bilinear estimates apply to the three deterministic
vectors \(\mathbf1/\sqrt n,u,Z_1\), whose norms are bounded by the retained
second-moment bound. Their transformed Gram entries are close to
\(v^\top Av'\) for the corresponding pair \(v,v'\); the extra vector \(Y_A(W_0u)\) has mixed entries close to
\(\chi v^\top Au\) and squared norm close to
\(\chi+\chi^2u^\top Au\). These are Gram comparisons under the fresh
\(W_0\) law. At the limiting scalar equality law with \(A=V\) and
\(\sigma=0\), the whitened finite-rank test has
determinant
\[
 (1-b+2g_1)(1-b+g_1)^2>0,
 \qquad
 g_1=\mathbb E_{N(s,s)}[\tanh^2Y\,\operatorname{sech}^2Y].
\]
The rank perturbation has at most one possibly nonpositive direction; the
positive determinant excludes that direction, exactly as in
\cite[Lemma~5.5]{AcceptedGap}. Compactness gives a uniform margin. For $A$
close to $V$, choose $\eta$ and its normalized Frobenius tolerance so that
$\|A\|\langle A\rangle\le1-\kappa$ on the retained range. To include
data-dependent \(A\), keep \(y\) fixed and round the entries exceeding a fixed sup-norm tolerance on
their exceptional coordinate set. That set has fraction at most
$\epsilon_A^2/\zeta^2$, and the logarithm of the number of choices is at most
\[
 n\{H(\epsilon_A^2/\zeta^2)
             +(\epsilon_A^2/\zeta^2)\log(C/\zeta)+o(1)\}.
\]
Choose \(\zeta\) after fixing the vector-norm bounds, Gram tolerance, and
fixed-diagonal exponent, then \(\epsilon_A\) so that this counting rate is
below that exponent. For this fixed \(y\), the approximants are deterministic,
so a union bound makes the conditional matrix comparison simultaneous in
\(A\). The omitted scalar mass has norm at most
\((1+\eta)|\chi-b|\le(1+\eta)n^{-1/2}\|A-V\|_{\rm F}\) and is included in
the reserved margin. The probability bound is uniform in the retained
\(y\)-data, as required inside the scalar integral. Enlarging the retained ranges slightly before the
local-volume step preserves all these margins.

The preceding calculations are under the usual size-biased planted law. The
critical zero-field partition comparison in
Proposition~\ref{cap:static-comparison} gives $\log Z-\log\mathbb EZ=o(n)$ in
probability. Thus negative exponential rates transfer to typical original
disorder: restrict to $Z/\mathbb EZ\ge e^{-\delta n}$, choose $\delta$ below
the rate, and use Markov's inequality. A finite time mesh and the Brownian
maximal bound supply the whole path. Finally, $F_h$ is one-Lipschitz in its
field and its stability matrix does not use $h$; coefficient and diagonal
changes are uniformly continuous under the norm bounds. Reserving slack in all
thresholds proves the neighborhood and perturbation assertions. \end{proof}

\subsection{Noise diagnostics at a fixed small scale}

We next work at $j=1$ with the full GOE matrix $W$. A spectral Gaussian means
$g=O\Sigma^{1/2}\xi$, where $W=O\Lambda O^{\mathsf T}$, $\xi$ is standard
Gaussian independent of $O$, and $\Sigma$ is positive semidefinite and fixed
after conditioning on the eigenvalues. We allow a fixed smooth compact family
of such covariances with $\|\Sigma\|\le Ct$. Write the field as
\begin{equation}\label{cap:boundary-field}
 h=g+h_1,\qquad |h_1|_n\le Ct.
\end{equation}
The shift may be selected after seeing the frame and noise in every frame
estimate below. A lower bound for a shifted Gaussian projection has a
different quantifier, specified separately in Lemma~\ref{cap:boundary-noise}.

Let $\Pi_d$ denote the projections onto dyadic bands for the gap $2-W$, with rank at most
$Cnd^{3/2}$. At a bottom label $a$, combine all smaller gaps into one band
with that label. Separately, write $P_t^W=\mathbf1_{\{2-W\le t\}}$ for
the cumulative projection onto gaps at most $t$; thus its cap-width parameter
is $\sqrt t$. If a vector class obeys $|\Pi_d v|_n\le b_d$, put
\begin{equation}\label{cap:band-complexity}
 \mathscr L(v)=\sum_{d\ge a}b_d d^{3/4}.
\end{equation}
Bounds for the individual band sizes can be chosen on arbitrarily fine fixed
grids. A finite collection of known spectral-coordinate vectors, including the
filtered noise, is included in the Gaussian comparison but is not charged this
complexity. Its number is fixed independently of $n$; smooth scalar spectral
filters can be netted before $n$ increases.

We record the two frame estimates used repeatedly.  Up to logarithmic
factors in the band count, for a vector of normalized size at most $B$,
\begin{equation}\label{cap:sparse-restriction}
 \sup_{|E|\le fn}|\mathbf1_Ev|_n
 \le C\{B\sqrt{f\log(2/f)}+\mathscr L(v)\}.
\end{equation}
Empirical averages of smooth clipped functions of finitely many vectors,
normalized to have size at most one, differ from the corresponding centered
Gaussian averages with the same Gram matrix by
\begin{equation}\label{cap:frame-comparison}
 C_\psi\sum_v\mathscr L(v).
\end{equation}
Here the complexities are those of the normalized vectors and $C_\psi$ is
controlled by the first two derivatives of the clipped function. Both
estimates permit an additional fixed slack and a failure allowance
$C\sqrt{\mathfrak h/n}$, multiplied by the indicated size and derivative
bounds, at failure rate $e^{-c\mathfrak h}$. We use $\mathfrak h$ of order
$nt^{3/2}$.

For completeness, \eqref{cap:sparse-restriction} follows by testing each band
against a sparse coordinate unit vector. Nets of its coefficient sphere and of
the chosen rows have logarithmic size $Cnd^{3/2}+Cnf\log(2/f)$, respectively.
The Gaussian representation of a Haar unit vector bounds its scalar products
at the square root of this entropy divided by $n$. Sum the band bounds and use
Cauchy--Schwarz for the total-size term. To prove
\eqref{cap:frame-comparison}, first fix \(k\) spectral coefficient columns
\(C=(c^1,\ldots,c^k)\), with \(k\) fixed and \(|c^j|_n\le1\). Put
\[
 Q=C^\top C/n,\qquad
 F_\psi(O,C)=n^{-1}\sum_i\psi((OC)_i).
\]
For \(\omega\) uniform on \(S^{n-1}\), rotation invariance gives the exact
one-row expectation
\[
 \E_OF_\psi(O,C)
 =\E_\omega\psi\bigl(\sqrt n\,Q^{1/2}(\omega_1,\ldots,\omega_k)^\top\bigr)
 =\E\psi(G_Q)+o_n(1),\qquad G_Q\sim N(0,Q).
\]
The formula is interpreted on the support of a possibly singular \(Q\).
Thus the Gaussian uses the actual Gram matrix of the chosen columns.
For the usual Frobenius metric on the orthogonal group,
\[
 \operatorname{Lip}_O F_\psi
 \le C_\psi\|C\|_{\rm op}/\sqrt n=O(C_\psi),\qquad
 \operatorname{Lip}_O(F_\psi(\cdot,C)-F_\psi(\cdot,C'))
 \le C_\psi\|C-C'\|_{\rm F}/\sqrt n.
\]
The first estimate follows by differentiating the row average and using
\(\|\nabla\psi(OC)\|_{\rm F}\le C_\psi\sqrt n\); the second uses the
bounded second derivative. Haar concentration has exponent proportional to
\(n\) at this Lipschitz scale. Chaining over the band balls therefore costs
\(C_\psi\sum b_d\sqrt{\operatorname{rank}\Pi_d/n}\), and a failure budget
\(\mathfrak h\) costs \(C_\psi\sqrt{\mathfrak h/n}\).
This proves \eqref{cap:frame-comparison}. The event is uniform over the
coefficient balls, so vectors may be selected from those balls after the
frame is observed.

We intersect the full-frame event with the following residual-rotation
event. For this calculation fix the eigenvalues \(\Lambda\), the low
physical frame \(O_0\), and a high spectral multiplier \(D\) at a fixed
scalar net point. Choose an isometry \(H\) onto the complement of the low
frame and write the remaining columns as \(O_>=HU\), where \(U\) is Haar
in dimension \(m\). Put \(M=HUDU^\top H^\top\) and
\(P_\perp=HH^\top\). Its exact conditional mean is
\[
 \E[M\mid\Lambda,O_0,D]=\mu P_\perp,\qquad \mu=\Tr D/m.
\]
For a finite net of vectors \(b,t\) fixed under this conditioning, let
\(h_{\rm net}\) be its logarithmic cardinality and \(h_{\rm tail}>0\).
The Haar quadratic-form bound and polarization give, with
\(\mathfrak H=C(h_{\rm net}+h_{\rm tail})\), outside conditional probability
\(2e^{-h_{\rm tail}}\),
\begin{equation}\label{cap:conditional-high-test}
 \begin{split}
 |\langle b,Mt\rangle-\mu\langle P_\perp b,P_\perp t\rangle|
 \le C|b|_n|t|_n\left(
 \frac{\|D\|_{\rm F}}{\sqrt m}\sqrt{\frac{\mathfrak H}{m}}
       +\|D\|\frac{\mathfrak H}{m}\right).
 \end{split}
\end{equation}
The probability is over the remaining \(U\); coefficient and sparse nets
are included in \(h_{\rm net}\). This is the conditional estimate used for
the high multipliers below. The low coordinate profiles remain fixed.

We specify the finite families included in the noise event. Scalar meshes
cover the covariance parameters, \(r\), the small values \(s=q+\sqrt t\),
and the derivative shift \(0\le\ell\le Cs\), with cutoffs
\(a=s^{7/5}\) and \(u=s^{1/10}\).
The uncharged spectral columns are the noise and the finitely many scalar
resolvent filters used in the root and derivative bases, including
\(r^{-1}\mathbf1_{\{d>a\}}(r+r^{-1}-\Lambda)^{-1}\Sigma^{1/2}\xi\).
The normalized root, derivative, base, correction, and band-test vectors
range over their charged coefficient balls. Sparse tests cover the root and
derivative tails, their union with a band-test tail, and the exceptional
\(A-V\) coordinates. The clipped recipe types are the root and derivative
sources against one band test, the scalar functions \(m^2,m^4,m^6,V'_0(y)w\),
their low-input versions, and the differentiated functions in the single
second-order Taylor expansion below. The number of these types and the
number of columns in each average are fixed before \(n\); the possible
directions within a type are paid for by the coefficient and sparse nets.
For each type the event is uniform over every band ball whose radii lie on
the chosen fixed grids; the later localization bounds select those radii.

\begin{lemma}[Uniform shifts and conditional noise]\label{cap:boundary-noise}
At every fixed sufficiently small $t>0$, the spectral band, sparse-row,
and clipped-average diagnostics just described can be imposed on an
event depending only on $(W,g)$ and the fixed parameter family.
They hold there uniformly over all shifts in \eqref{cap:boundary-field}
and the coefficient balls, sparse supports, and recipe and filter types
specified above. Their conditional failure
probability tends to zero for typical $W$; with fixed buffered tolerances
it is at most $e^{-c_tn}$ after decreasing $c_t>0$.  The event is independent
of any square-density test function.

Suppose additionally that the covariance of $g$ is at least $ct$ on a
spectral band contained in $\{2-W\le t\}$ of rank at least $cnt^{3/2}$.
Then, for every fixed center $v\in\mathbb R^n$,
\begin{equation}\label{cap:small-ball}
 \mathbb P_g\bigl\{|P_t^W(g+v)|_n<c't^{5/4}\mid W\bigr\}
 \le e^{-c''nt^{3/2}}.
\end{equation}
The constants are uniform in the center.  This is a probability bound
for each center, not a simultaneous event over all centers.
In particular, both types of estimate apply after averaging over any
spin law that is fixed before the fresh Gaussian noise is drawn.
\end{lemma}

\begin{proof}
There are two different uniformity statements here. The frame estimates
hold on one event for all coefficient vectors in the specified balls, so
they admit shifts selected after the frame and noise are seen. The
projection lower bound instead holds in probability for each fixed center,
with a constant independent of that center. A spin sampled before the
fresh noise can therefore be averaged into either statement. There is no
simultaneous lower bound over every possible center.

The frame estimates are simultaneous coefficient-ball estimates. The shift
contributes at most $Ct$ to each unfiltered band norm; after a high-band
inverse it contributes at most $Ct/d$. Thus every shift in the stated ball is
already included in their deterministic domains. The Gaussian band norms and
the norms of their finite smooth filtered families have exponential tails at
the specified fixed scales. If \(p(W)\) is the conditional noise failure
probability and \(\E_Wp(W)\le e^{-2c_tn}\), Markov gives
\(\Pr_W\{p(W)>e^{-c_tn}\}\le e^{-c_tn}\). This gives the asserted disorder
event and conditional bounds, without any reference to a test function.
The same argument integrated over a fixed time interval gives
the integrated conditional statement, when that is all that is needed.

For \eqref{cap:small-ball}, restrict to the nondegenerate band of rank $r_t\ge
cnt^{3/2}$. The Gaussian integral of a Euclidean ball is maximized when its
center is the mean; alternatively, exponential Markov applied to
$\exp(-\lambda\|g+v\|^2/t)$ gives a product bounded by
$(1+c\lambda)^{-r_t/2}$, since the noncentral exponential factors are at most
one. Taking the ball radius squared to be a sufficiently small multiple of
$tr_t$ proves the claim. Finally, condition on each spin and then average this
bound over its law. No union bound over spin configurations is involved, and
no initial density bound is used. \end{proof}

\subsection{The small-field TAP expansion}

\begin{proposition}[Boundary expansion]\label{cap:tap-expansion}
Assume the preceding spectral and noise diagnostics, the zero-field
approximate-root exclusion in Lemma~\ref{cap:endpoint-stability}, the field
representation \eqref{cap:boundary-field}, and
$|P_t^Wh|_n\ge ct^{5/4}$. The test \eqref{cap:stability-test} holds with
positive fixed slack; in particular $F_h$ has a unique root $y_*(h)$.
At all roots of $F_h$,
and with arbitrarily small additional errors at sufficiently accurate
approximate roots, set
\[
 z=r+r^{-1},\qquad R=(z-W)^{-1}.
\]
For sufficiently small fixed $t$ and then sufficiently large $n$,
\begin{equation}\label{cap:tap-leading}
 \begin{gathered}
 q\ge c\sqrt t,\qquad |y-r^{-1}Rh|_n=o(\sqrt q),\\
 K_{\mathrm{TAP}}
 =R-\frac{4q\,Ryy^{\mathsf T}R/n}
                 {1+4q\langle yRy\rangle}+o(q^{-2}).
 \end{gathered}
\end{equation}
Here $K_{\mathrm{TAP}}=\partial_h\tanh y_*(h)$, and the last error
is in operator norm. The errors tend to zero as the boundary scales
decrease, after which all dimensional errors and perturbation tolerances
are chosen.

If $q$ is comparable to $\sqrt t$, define the TAP potential per site by
\[
 \mathcal T(h)=\langle hm\rangle+\tfrac12\langle mWm\rangle
                  -\langle I(m)\rangle+\tfrac14r^2,
 \quad
 I(v)=\tfrac12(1+v)\log(1+v)+\tfrac12(1-v)\log(1-v).
\]
With $\kappa_0=1/3$, the value expansion is
\begin{equation}\label{cap:tap-value}
 \mathcal T(h)-\tfrac12\langle hRh\rangle
 =\frac{r^2}{4}+\frac{zq}{2}-\langle I(m)\rangle
                    -\frac{\langle eRe\rangle}{2r^2}
 =\frac14+\kappa_0q^3+o(q^3),
 \qquad e=\tanh y-ry.
\end{equation}
All assertions persist to any prescribed small macroscopic error under
sufficiently small changes of the interaction coefficient, operator-norm
perturbations, and normalized field displacements, with the TAP potential
using the actual interaction and $jr^2/4$ at its actual root.
\end{proposition}

\begin{proof}
We first control the nonlinear source on the coordinates where a root
is large, uniformly over every approximate root in the specified region.
This gives initial localization bounds for the root and its derivative
tests. A second gain is then necessary: an error of order $q^2$ would still
be as large as the spectral separation. The later high-rotation calculation
provides that gain before we invert on the low bands.
Lemma~\ref{cap:endpoint-stability} at zero, applied with the small field
as an additional residual, first makes $q$ as small as required.
The field equation gives $|y|_n\le C(\sqrt q+\sqrt t)$.
Put
\begin{equation}\label{cap:clipping-scales}
 s=q+\sqrt t,\qquad a=s^{7/5},\qquad
 S=\{2-W\le a\},\qquad u=s^{1/10},\qquad
 P_0=\mathbf1_{\{2-W\le u\}}.
\end{equation}
We only invert on $S^\perp$ until the lower bound for $q$ proves separation.
Dyadic bands there have gaps between $a$ and a constant. In all estimates
below $o_n(1)$ errors at fixed scales are chosen after the smallest scale and
may be below any specified power of $s$.

The derivative estimates below test vectors $w$ normalized by $|w|_n\le1$.
For a diagonal $A$ allowed by the stability test, put $v=Aw$ and
$\dot r=-2\langle mv\rangle$. Equation \eqref{cap:derivative-equation}
will specify the equation these tests satisfy; the present notation lets
us account for their clipping errors before that calculation.
The scalar variance profile satisfies
$V'_0(z)=-2\tanh z\,\operatorname{sech}^2z$. The cutoffs below are applied
to recipes built from this profile, including $V'_0(y)w$.

\paragraph{Clipping convention and its power budget.}
Write $T_s=s^{-\kappa}$, where $\kappa>0$ is chosen after the finite
clip-loss constant specified below.  Use a smooth cutoff that is one on
$[-1,1]$, zero off $[-2,2]$, and has fixed bounded derivatives.  Its
arguments are $y/(\sqrt sT_s)$, $w/T_s$, and a normalized linear test
divided by $T_s$.  The normalized root source
$e/s^{3/2}$ and derivative source
$[(q-m^2)w-\dot r\,y]/s$ have degree-three envelopes on this region;
with a linear test they have degree-four envelopes.  The normalized
scalar recipes for $m^2,m^4,m^6$, and $V'_0(y)w$ have envelopes of
degrees $2,4,6$, and $2$, respectively.  In these statements ``degree''
means that every derivative through the order used is bounded by a
constant times that power of $T_s$.  The frame comparison uses at most
two derivatives, and its differentiated Taylor tests use at most three.
Cutoff differentiation contributes a nonpositive power of $T_s$.
Thus all the needed envelopes are polynomial in $T_s$.

The proof uses a fixed finite sequence of operations: band projection and
dyadic summation, sparse inverses of norm at most two, two complexity
absorptions, replacement of high inputs by \(P_0\) inputs, and one
second-order Taylor expansion. The scalar moment and value estimates use the
same recipes and one further product or contraction. The polynomial
derivative bounds above and the logarithmic sizes of the Gaussian and Haar
nets therefore give each displayed error the form
\(Cs^\alpha T_s^{C_i}\), with a finite \(C_i\) independent of
\(\kappa,s,n\). The powers from \(a^{-1}\) and \(u^{-1}\) are displayed
as powers of \(s\); they are not included in \(C_i\).

There is a corresponding finite bound on the clipping exponents in the
absorption coefficients. In an inequality
\(X\le Cs^\alpha T_s^c+Cs^\theta T_s^dX\) with \(\theta>0\),
the second term is absorbed once \(\kappa d<\theta\) and the upper scale is
small. Let \(C_{\rm clip}\) dominate all the finitely many exponents
produced by these operations, including those \(d\). Choose
\begin{equation}\label{cap:clip-choice}
 0<\kappa<\frac1{1000(1+C_{\mathrm{clip}})}.
\end{equation}
This choice precedes the upper boundary scale. Below, \(T_s^C\) denotes one
of these fixed losses. The smallest raw gain in the displayed estimates is
\(1/20\), so \eqref{cap:clip-choice} leaves a positive power in each
absorption and in the final errors.

The natural-size term in every sparse restriction is kept as
$CB\sqrt{f\log(2/f)}$, with no $T_s^C$ multiplier. Only its complexity term
incurs clipping losses. Thus threshold absorption uses
$\sqrt{\log(2/f)}/T_s\to0$ and is not circular. First separate fractions below
an arbitrarily high power of $s$, whose contributions already lie below the
target errors; otherwise $\log(2/f)=O(|\log s|)$. The finite remaining
logarithms fit into a further fixed power of $T_s$. The failure allowance
$\sqrt{\mathfrak h/n}=O(t^{3/4})\le O(s^{3/2})$ also lies below the errors for
which it is used.

\paragraph{The first localization of the root.}
We first localize the root strongly enough to make the nonlinear source
small on high spectral bands. This first estimate is intentionally weaker
than the final low-band cancellation: it controls the tail set and makes a
second approximation possible. The low-band root component and the filtered
shift form the charged part of the base vector; the known Gaussian noise is
included in the frame comparison without paying a vector-class complexity.

The exact high-band equation is
\begin{equation}\label{cap:root-band-equation}
 y=b_y+R'_{\mathrm{hi}}e,\qquad
 b_y=P_Sy+r^{-1}R_{\mathrm{hi}}h,\qquad
 R'=r^{-2}R-r^{-1}I.
\end{equation}
Both $y$ and $b_y$ have normalized size $O(\sqrt s)$.  The uncharged
noise term has band norm $O(s d^{-1/4})$.  The low-band vector and the
filtered shift give
\begin{equation}\label{cap:base-root-complexity}
 \mathscr L(b_y)
 \le C\{\sqrt s\,a^{3/4}+s^2a^{-1/4}\}
 =C\{s^{31/20}+s^{33/20}\},
\end{equation}
up to logarithms.  Let $L_y$ be this complexity plus the complexities of
the correction bands in \eqref{cap:root-band-equation}.  Initially
$L_y=O(\sqrt s)$ is sufficient.

Choose a smooth scalar cutoff $0\le\chi\le1$ equal to one on $[-1,1]$ and
zero outside $[-2,2]$, and put
$e_0(y)=\chi(y/(\sqrt sT_s))(\tanh y-ry)$ coordinatewise.
Then $|e_0|\le Cs^{3/2}T_s^3$ coordinatewise and $|e'_0|\le CsT_s^2$.
On the transition and tail coordinates, $\tanh y/y\le r$ for sufficiently
small $s$, so
\[
 e-e_0=-D_Ey,\qquad 0\le D_E\le rI,
 \qquad E\subseteq\{|y|>\sqrt sT_s\}.
\]
The sparse restriction estimate, with its first term absorbed against $\sqrt
sT_s\sqrt f$, gives $\sqrt f\le Cs^{-1/2}T_s^{-1}L_y$, apart from the already
reserved tiny errors. Testing the clipped source on an arbitrary normalized
unit vector in a band gives
\begin{equation}\label{cap:clipped-root-source}
 |\Pi_de_0|_n\le CT_s^C
                 \{s^{3/2}d^{3/4}+sL_y\}.
\end{equation}
Here is the cancellation in this estimate. Gaussian integration by parts
predicts its linear projection with coefficient $\mathbb E e'_0(Y)$, where
$\operatorname{Var}Y=|y|_n^2$. The identity $\langle q-m^2\rangle=0$ and the
frame comparison bound the corresponding unclipped coefficient by
\(CT_s^C\sqrt sL_y\). For the cutoff error, the threshold estimate above
and sparse restriction give
\[
 \sqrt f\le Cs^{-1/2}T_s^{-1}L_y,\qquad
 |\mathbf1_Ey|_n\le CT_s^CL_y .
\]
Since the coarse input is \(L_y=O(\sqrt s)\),
\[
 |\mathbf1_Ey|_n^2+sf
 \le CT_s^CL_y^2\le C\sqrt s\,T_s^CL_y .
\]
This bounds the change in \(\mathbb E e'_0(Y)\). Multiplication by
\(|\mathbb E[YU]|\le C\sqrt s\) for the normalized Gaussian band test \(U\)
gives the \(sL_y\) term in
\eqref{cap:clipped-root-source}. The centered part has band-complexity cost
\(Cs^{3/2}T_s^Cd^{3/4}\). Clipping a linear test
first and then removing its clip costs its complexity times the size of the
bounded source, again by \eqref{cap:sparse-restriction}.

Summing \eqref{cap:clipped-root-source} with weight $d^{-1/2}$ yields
\[
 |R_{\mathrm{hi}}^{1/2}e_0|_n+|e_0|_n
 \le CT_s^C(s^{3/2}+sa^{-1/2}L_y).
\]
The tail equation is coercive: on the high bands,
$I+D_E^{1/2}R'_{\mathrm{hi}}D_E^{1/2}$ dominates a fixed positive multiple of
$I+D_E^{1/2}R_{\mathrm{hi}}D_E^{1/2}$. Indeed, the initial bound $|y|_n\le
C\sqrt s$ already gives $f=|E|/n\le CT_s^{-2}$. The uniform sparse restriction
bound on $S$ therefore implies
\[
 \|P_S\mathbf1_E\|
 \le C\{a^{3/4}+\sqrt{f\log(2/f)}\}+o_n(1)=o(1),
\]
before any improved tail estimate. Put $X_E=D_E^{1/2}R_{\mathrm{hi}}D_E^{1/2}$
and $Q_E=D_E^{1/2}P_SD_E^{1/2}$. Since $D_E$ is supported on $E$, $\|Q_E\|\le
r\|P_S\mathbf1_E\|^2=o(1)$, and $R_{\mathrm{hi}}\succeq cP_{\mathrm{hi}}$
gives $X_E\succeq cD_E-o(1)I$. Thus, for a sufficiently small fixed
$\theta>0$,
\[
 \begin{aligned}
 &I+D_E^{1/2}R'_{\mathrm{hi}}D_E^{1/2}-\theta X_E\mathstrut\\
 &\quad=I-r^{-1}D_E+r^{-1}Q_E+(r^{-2}-\theta)X_E\\
 &\quad\succeq I+\{c(r^{-2}-\theta)-r^{-1}\}D_E-o(1)I\\
 &\quad\succeq c'I,
 \end{aligned}
\]
where \(c'>0\) is fixed because \(D_E\le rI\) and \(r\to1\). This
domination uses only the coarse tail fraction.

To apply it, put \(u_E=D_E^{1/2}y\). Substituting \(e=e_0-D_Ey\) into
\eqref{cap:root-band-equation} gives the exact tail equation
\[
 (I+D_E^{1/2}R'_{\rm hi}D_E^{1/2})u_E
 =D_E^{1/2}b_y+D_E^{1/2}R'_{\rm hi}e_0 .
\]
Test against \(u_E\) and use the lower bound by a fixed multiple of
\(I+X_E\). Decomposing \(R'_{\rm hi}=r^{-2}R_{\rm hi}-r^{-1}P_{\rm hi}\)
and applying Cauchy--Schwarz in that energy norm bounds the right side by
the dual sizes \(|D_E^{1/2}b_y|_n\), \(|e_0|_n\), and
\(|R_{\rm hi}^{1/2}e_0|_n\). Since \(D_E\le rI\), the resulting estimate is
\begin{equation}\label{cap:root-tail-energy}
 |D_Ey|_n+|R_{\mathrm{hi}}^{1/2}D_Ey|_n
 \le C\{|D_E^{1/2}b_y|_n+|e_0|_n
                         +|R_{\mathrm{hi}}^{1/2}e_0|_n\}.
\end{equation}
The first term is $o(1)L_y+C\mathscr L(b_y)$ by sparse restriction. The high
correction complexity is bounded by the right side of
\eqref{cap:root-tail-energy}; the coefficient $sa^{-1/2}=s^{3/10}$ absorbs
after clipping losses. We conclude
\begin{equation}\label{cap:first-root-localization}
 \begin{gathered}
 L_y\le Cs^{3/2}T_s^C,\qquad f\le Cs^2T_s^C,\\
 |y-b_y|_n\le CT_s^C
       \{s^{3/2}a^{-1/4}+s^{5/2}a^{-1}\}
       \le Cs^{11/10}T_s^C.
 \end{gathered}
\end{equation}
For the last bound, the sparse part is tested bandwise using
$\|\Pi_d\mathbf1_E\|\le CT_s^C(d^{3/4}+s)$; the nonsparse part uses
\eqref{cap:clipped-root-source}. This also shows $|e|_n\le Cs^{3/2}T_s^C$.

\paragraph{Localization of derivative tests.}
Normalize $|w|_n\le1$, put $v=Aw$ and
$\dot r=-2\langle mv\rangle$, and consider
\begin{equation}\label{cap:derivative-equation}
 w+(rI-W)v+\dot r\,m=h_v-\ell v,
 \qquad |h_v|_n\le Cs^2,\quad 0\le\ell\le Cs.
\end{equation}
This includes possible nonpositive stability eigenvalues, after a scalar
shift, and the response to bounded sources. Set $R_\ell=(z+\ell-W)^{-1}$,
$R'_\ell=r^{-2}R_\ell-r^{-1}I$, and $d_v=(A-rI)w-\dot r\,y$. Exact elimination
gives
\begin{equation}\label{cap:derivative-band-equation}
 \begin{split}
 w&=b_w+R'_{\ell,\mathrm{hi}}d_v,\\
 b_w&=P_Sw+r^{-1}R_{\ell,\mathrm{hi}}h_v\\
 &\quad+\dot r\{R_{\ell,\mathrm{hi}}((r^{-2}-1)y-e/r)
                                    -r^{-1}P_{\mathrm{hi}}y\}.
 \end{split}
\end{equation}
Using the noise band sizes, \eqref{cap:first-root-localization}, and
$|\dot r|\le C\sqrt s$ gives
\[
 \mathscr L(b_w)\le CT_s^C
           (a^{3/4}+s^{3/2}+s^2a^{-1/4}),
 \qquad |P_{\mathrm{hi}}b_w|_n
        \le CT_s^C(s+s^2/a).
\]
The first bound starts at $s^{21/20}$; this small gain will be retained.
Let $L_w$ include the complexity of the correction bands.

Split $d_v=d_0+D'_Ew+e_A$, where $d_0$ uses the bounded clipped coefficient
$q-m^2$ and retains $-\dot r\,y$. Here $D'_E$ is bounded and supported on the
tail set; $e_A$ is an arbitrarily small normalized error. To obtain this
decomposition for all allowed $A$, fix a diagonal sup-norm tolerance
$\delta_A$, split $A-V$ outside that tolerance, and use
\[
 \frac{|\{i:|A_{ii}-V_{ii}|>\delta_A\}|}{n}
 \le\epsilon_A^2/\delta_A^2.
\]
At fixed bottom scale choose $\delta_A$ and then $\epsilon_A$ so that the
first part is below every requested error and the added fraction is below
every requested power of $s$. This choice is made before $n$, and the frame
tests remain simultaneous over the added sparse set.

Uniformly on the high bands,
\begin{equation}\label{cap:clipped-derivative-source}
 |\Pi_dd_0|_n\le CT_s^C
        \{sd^{3/4}+sL_w+\sqrt sL_y\}.
\end{equation}
For jointly centered Gaussian surrogates $Y,G,U$, Gaussian integration
by parts gives the exact identity
\[
 \mathbb E[\operatorname{sech}^2(Y)GU]
 =\mathbb E\operatorname{sech}^2(Y)\,\mathbb E[GU]
   +\mathbb E[(\operatorname{sech}^2)'(Y)G]\,\mathbb E[YU].
\]
Centering by $r$ removes the first term; $\dot r$ removes the second. The
frame comparison for the latter scalar coefficient has natural size $\sqrt s$,
and its unclipping cost is bounded by $C|\mathbf1_Ey|_n|\mathbf1_Ew|_n$. This
proves \eqref{cap:clipped-derivative-source}; the diagonal error is handled by
the preceding split. Removing clips on its linear factors uses sparse
Cauchy--Schwarz, so it charges the sum of their complexities, not their full
norms.

The compression of \(M=R'_{\ell,\mathrm{hi}}\) to the enlarged sparse set
has operator norm \(o(1)\). Its spectral mean is \(O(\sqrt a+\sqrt s)+o_n(1)\), its
Frobenius norm divided by $\sqrt n$ is at most $Ca^{-1/4}$, and its operator
norm is at most \(C/a\). Put
\(r_f=f\log(2/f)+\mathfrak h/n\). The Haar quadratic-form bound and the
sparse nets give
\[
 \|\mathbf1_E M\mathbf1_E\|
 \le \frac{|\Tr M|}{n}
      +C a^{-1/4}\sqrt{r_f}+C a^{-1}r_f+o_n(1).
\]
Since \(f\le Cs^2T_s^C\), \(\mathfrak h/n=O(t^{3/2})=O(s^3)\), and
\(a=s^{7/5}\), the two deviation terms have raw powers at least
\(s^{13/20}\) and \(s^{3/5}\). The clip choice makes them tend to zero.
Restricting \eqref{cap:derivative-band-equation} to the sparse coordinates
gives
\[
 (I-\mathbf1_E M\mathbf1_E D'_E)\mathbf1_Ew
 =\mathbf1_Eb_w+\mathbf1_E M(d_0+e_A).
\]
Because \(D'_E\) is bounded, the inverse on the left has norm at most two
after the upper scale is decreased. The \(e_A\) term is below the reserved
error. Bandwise testing of the other term gives
\begin{equation}\label{cap:derivative-tail-energy}
 \begin{split}
 |\mathbf1_Ew|_n\le C\bigg[|\mathbf1_Eb_w|_n+
 T_s^C\sum_{d\ge a}d^{-1}(d^{3/4}+\sqrt f)
             (sd^{3/4}+sL_w+\sqrt sL_y)\bigg].
 \end{split}
\end{equation}
The factor multiplying $sL_w$ sums to $O(a^{-1/4}+sT_s^Ca^{-1})$, before the
displayed clipping loss. Multiplication by $s$ thus gives a positive power;
the sparse-source complexity costs only $C(1+sa^{-1/4})$, apart from clips and
logarithms. Together with the preceding bound for $\mathscr L(b_w)$ this
absorbs the $L_w$ terms and proves
\begin{equation}\label{cap:first-derivative-localization}
 L_w\le CsT_s^C,\qquad
 |P_{\mathrm{hi}}w|_n
 \le CT_s^C(sa^{-1/4}+s^2/a)
 \le Cs^{3/5}T_s^C.
\end{equation}

\paragraph{A second gain from the independent high rotation.}
The first root and derivative estimates have not yet beaten the spectral
separation $s^2$. We now separate each vector into a part with smaller complexity
and a structured high-frequency correction. The middle bands provide a
power gain by summation. On the remaining high bands, replacing each source
by a recipe of its low-frequency inputs leaves the high frame available for
a signed estimate. Full operator-norm bounds would lose this gain, so the
contractions are estimated before taking absolute operator norms.

  We now prove a stronger decomposition,
with some fixed $\zeta>0$; one may take $\zeta=1/100$ after
\eqref{cap:clip-choice}:
\begin{equation}\label{cap:refined-decomposition}
 \begin{array}{ll}
 y=B_y+R'_{>u}t_y,&
   |t_y|_n\le Cs^{3/2}T_s^C,\quad
       \mathscr L(B_y)\le Cs^{3/2+\zeta},\\
 w=B_w+R'_{\ell,>u}t_w,&
   |t_w|_n\le CsT_s^C,\quad
       \mathscr L(B_w)\le Cs^{1+\zeta}.
 \end{array}
\end{equation}
As before, the filtered noise is included but not charged in $\mathscr
L(B_y)$. To specify the recipes, put $Y_0=P_0y$, $W_0=P_0w$ and let $\chi$ be
the preceding scalar cutoff. We may take
\begin{equation}\label{cap:low-input-recipes}
 \begin{split}
 t_y&=\chi\left(\frac{Y_0}{\sqrt sT_s}\right)
                         (\tanh Y_0-rY_0),\\
 t_w&=\chi\left(\frac{Y_0}{\sqrt sT_s}\right)
       \chi\left(\frac{W_0}{T_s}\right)
       \{(\operatorname{sech}^2Y_0-r)W_0-\dot r\,Y_0\}.
 \end{split}
\end{equation}
All products here are coordinatewise. Small operator and sparse $A-V$ errors
are assigned to $B_w$, using the tail equations below. Net the scalar
parameters on their bounded ranges and let the low inputs vary over all
vectors of the stated sizes in $P_0$. Once these coefficients and the low
columns have been fixed, the recipes are determined. The high columns are
still free to rotate, which is the independence used in the next signed
estimate.

First sum the source bounds on $a\le d\le u$.  Both the weighted
complexity and the norm restricted to a fraction
$f\le s^2T_s^C$ gain, relative to the source size, the factor
\begin{equation}\label{cap:middle-band-gain}
 T_s^C(u^{1/2}+sa^{-1/4}+s^2a^{-1})
       =T_s^C(s^{1/20}+s^{13/20}+s^{3/5}).
\end{equation}
This follows by multiplying \eqref{cap:clipped-root-source} and
\eqref{cap:clipped-derivative-source} by $d^{-1}(d^{3/4}+\sqrt f)$ and summing
dyadically. The same bounds with only the complexity weight prove the claimed
gain in complexity. Above $u$ we replace the bounded sources by
\eqref{cap:low-input-recipes}. The necessary input estimates are
\begin{equation}\label{cap:high-input-bounds}
 \frac{|P_{>u}y|_n}{\sqrt s}
      \le CT_s^C(s^{19/40}+s^{3/5}),
 \qquad |P_{>u}w|_n\le CT_s^Cs^{3/5}.
\end{equation}
For the first, the filtered noise contributes $s u^{-1/4}$ before division by
$\sqrt s$, and \eqref{cap:first-root-localization} contributes $s^{11/10}$;
the second follows from \eqref{cap:first-derivative-localization}. The source
derivatives have natural sizes $s$ for the root source, $s$ in the $w$
variable for the derivative source, and $\sqrt s$ in its $y$ variable after
clipping $w$. Including the high inverse norm $C/u$, the replacement errors
are at most
\[
 CT_s^Cs^{15/8}\quad\hbox{for }y,\qquad
 CT_s^C(s^{11/8}+s^{3/2})\quad\hbox{for }w.
\]
These are smaller than the gained sizes in \eqref{cap:refined-decomposition}.
Clipping the linear $w$ input at this stage uses its first sparse
localization: the bounded coefficient has size $sT_s^C$, and the resulting
error after division by $u$ is $O(s^{19/10}T_s^C)$. Thus the refined joint
tail estimate is not being assumed in constructing these recipes.

We justify the signed estimate for the structured high terms. Conditional on
$P_0$ and its frame, the remaining rotation is Haar. For either high
multiplier $M=R'_{>u}$ or $R'_{\ell,>u}$,
\begin{equation}\label{cap:high-multiplier-sizes}
 \left|\frac1n\operatorname{Tr}M\right|
       \le C(\sqrt u+\sqrt s)+o_n(1),\qquad
 \frac{\|M\|_{\mathrm F}}{\sqrt n}\le Cu^{-1/4},
 \qquad \|M\|\le C/u.
\end{equation}
Net the low coefficient balls and the bounded scalar parameters before
substituting the actual, possibly frame-dependent coefficients. At polynomial
resolution in \(s\), their logarithmic cardinality is
\(O(nu^{3/2}|\log s|)\). Here the fixed complement has dimension
\(m=n-O(nu^{3/2})\asymp n\). Equation~\eqref{cap:conditional-high-test}
bounds the deviation from the exact conditional mean
\(\mu\langle P_\perp b,P_\perp t\rangle\), with \(\mu=\Tr D/m\), by
\[
 C\left(u^{-1/4}\sqrt{u^{3/2}|\log s|}
                  +u^{-1}u^{3/2}|\log s|\right)
 \le Cu^{1/2}T_s^C
\]
times the two vector sizes. The mean itself is bounded by
\(C(\sqrt u+\sqrt s)+o_n(1)\) from \eqref{cap:high-multiplier-sizes}.
The estimate is uniform over the nets and then over the full
classes by their derivative bounds. Adding sparse coordinate unit tests costs
$O(nf\log(2/f))$, which is smaller than this entropy. Applying this test to a
sparse output gives a factor $s^{1/20}T_s^C$ relative to the source norm.
Applying it to another low-input recipe gives the same factor relative to the
product of the two norms. These are the two forms of the gain used below. The
operator norm $u^{-1}$ is not charged a second time to a signed contraction:
it already appears in the second term of the quadratic-form estimate above.

The known high-noise column causes no conditioning problem here. It was
included in the finite-column frame comparison, where its spectral
coefficients are independent of the frame. For the present signed tests, first
replace each differentiated input by its $P_0$ version. Estimate the discarded
part in full norm, including the high-noise contribution $s^{19/40}$ in
\eqref{cap:high-input-bounds}. The power calculation below puts this error at
higher order. The remaining recipe no longer uses the high-noise image. The
retained recipes therefore depend only on the low frame. This distinction is
essential: we do not condition on the image of the high noise and then call
the remaining rotation unrestricted Haar.

It remains to improve the unclipped sparse sources. First separate the
exceptional \(A-V\) coordinates using the previously chosen diagonal
tolerances; their contribution is below the errors in this paragraph.
On the root tail \(E\), the exact restriction is
\[
 (I+\mathbf1_E R'_{\mathrm{hi}}\mathbf1_E D_E)\mathbf1_Ey
 =\mathbf1_Eb_y+\mathbf1_E R'_{\mathrm{hi}}e_0.
\]
The preceding sparse compression estimate includes \(\ell=0\), so the
inverse on the left has norm at most two. This is the unweighted sparse
inverse; the earlier energy estimate on \(D_Ey\) alone would not control
\(\mathbf1_Ey\). The middle-band gain and the structured sparse-output
estimate then give
\[
 |\mathbf1_Ey|_n
 \le C\sqrt s\sqrt{f\log(2/f)}+Cs^{3/2+1/20}T_s^C.
\]
The threshold \(\sqrt sT_s\sqrt f\le|\mathbf1_Ey|_n\) absorbs the first
term, with its natural-size coefficient unchanged. Hence
\[
 \sqrt f\le Cs^{1+1/20}T_s^C,\qquad
 |\mathbf1_Ey|_n\le Cs^{3/2+1/20}T_s^C.
\]
The derivative sparse equation uses this same root-coefficient tail,
together with the separately controlled \(A-V\) set. Its base restriction
is now bounded by
\(C\sqrt{f\log(2/f)}+Cs^{1+1/20}T_s^C\). The inverse of norm at most two
and the middle and structured output gains therefore give
\[
 |\mathbf1_Ew|_n
 \le C\sqrt{f\log(2/f)}+Cs^{1+1/20}T_s^C
 \le Cs^{1+1/20}T_s^C.
\]

For a sparse source of row norm \(B_E\), its contribution to the weighted
high-band complexity is at most \(B_E T_s^C\) times
\[
 \sum_{d\ge a}d^{-1}d^{3/4}(d^{3/4}+\sqrt f)
 \le C(1+\sqrt f\,a^{-1/4}),
\]
by the sparse band test and dyadic summation. Thus the improved row norms
give the raw bounds
\begin{equation}\label{cap:refined-raw-complexity}
 \mathscr L(B_y)\le Cs^{3/2+1/20}T_s^C,\qquad
 \mathscr L(B_w)\le Cs^{1+1/20}T_s^C.
\end{equation}
The middle-band and replacement errors have at least these powers.
The clip choice weakens \eqref{cap:refined-raw-complexity} to
\eqref{cap:refined-decomposition} with \(\zeta=1/100\).

Separate tail estimates are insufficient for the products used in a derivative
test. We therefore clip the root, derivative vector, and low-band test on the
union of their tail sets. The refined decomposition bounds all three vectors
on this same set. Their natural-size terms can be absorbed against the
clipping thresholds, leaving a positive power gain for every product discarded
when the clips are removed.

If $z_S$ is any normalized unit vector in $S$, then $\mathscr L(z_S)\le
Ca^{3/4}=Cs^{21/20}$. Let $E$ be the union of the tails of $y/\sqrt s$, $w$,
and $z_S$, each at threshold $T_s$, together with the sparse $A-V$ set. The
first bounds put its fraction below \(s^2T_s^C\). The raw bounds
\eqref{cap:refined-raw-complexity}, the structured sparse-output estimate,
and sparse restriction give
\begin{equation}\label{cap:joint-tail-inequalities}
 \begin{split}
 |\mathbf1_Ey|_n&\le
 C\sqrt s\sqrt{f\log(2/f)}+Cs^{3/2+1/20}T_s^C,\\
 |\mathbf1_Ew|_n+|\mathbf1_Ez_S|_n&\le
 C\sqrt{f\log(2/f)}+Cs^{1+1/20}T_s^C.
 \end{split}
\end{equation}
The threshold inequality, after removing the arbitrarily smaller
$A-V$ fraction, is
\[
 T_s\sqrt f\le
 |\mathbf1_Ey|_n/\sqrt s
        +|\mathbf1_Ew|_n+|\mathbf1_Ez_S|_n.
\]
The unmultiplied natural-size terms in
\eqref{cap:joint-tail-inequalities} absorb.  Therefore
\begin{equation}\label{cap:joint-tail-gain}
 \sqrt f\le Cs^{1+1/20}T_s^C,\qquad
 |\mathbf1_Ey|_n\le Cs^{3/2+1/20}T_s^C,\qquad
 |\mathbf1_Ew|_n+|\mathbf1_Ez_S|_n
                       \le Cs^{1+1/20}T_s^C.
\end{equation}
For an explicit permitted diagonal tolerance, split $A-V$ at $s^{20}$ and take
its normalized Frobenius tolerance at most $s^{40}$. The exceptional fraction
is then at most $s^{40}$; its contribution is smaller than all the displayed
errors. Across a fixed interval of boundary scales, use its smallest scale in
these choices.

\paragraph{The low-band cancellations.}
All parts needed for the low-band estimate are now available. The base
vectors have small enough complexity for comparison with Gaussian
averages; the linear structured corrections use the signed high-rotation
bound; and the quadratic Taylor remainder is smaller still. We list the
powers separately before removing the joint clips. After these errors have
been bounded, the Gaussian predictions cancel because the overlap and its
derivative were defined from the same root.


\begin{equation}\label{cap:low-source-cancellation}
 |P_Se|_n=o(s^{5/2}),\qquad |P_Sd_v|_n=o(s^2).
\end{equation}
We give the power ledger explicitly. Normalize the variables as $y/\sqrt
s,w,z_S$. The base complexities in \eqref{cap:refined-decomposition} are
$s^{1+\zeta}$, $s^{1+\zeta}$, and $s^{21/20}$. Applying the frame comparison
to the clipped averages therefore gains more than one power of $s$ relative to
the source sizes $s^{3/2}$ and $s$. The structured normalized perturbations
have source norm $sT_s^C$ and full norm at most $s/u$ times a clipping loss.
Their second-order Taylor remainder consequently has relative size
$s^{9/5}T_s^C$. For a linear term use the signed high contraction, giving
relative size $su^{1/2}T_s^C=s^{21/20}T_s^C$. Replacing its differentiated
recipe by a $P_0$-input recipe has relative error at most
$su^{-1}(s^{19/40}+s^{3/5})T_s^C$. The resulting exponents, before clip
losses, are
\begin{equation}\label{cap:boundary-power-ledger}
 \begin{array}{c|c}
 \text{operation}&\text{relative power of }s\\ \hline
 \text{base frame comparison}&1+\min(\zeta,1/20)\\
 \text{structured linear contraction}&21/20\\
 \text{second-order Taylor remainder}&9/5\\
 \text{high-input replacement}&11/8\quad\text{and}\quad3/2.
 \end{array}
\end{equation}
The same calculation controls changes of Gram entries in the Gaussian
predictions; the discarded high parts are included in the replacement error.
The differentiated recipes are exactly derivatives, through order two, of $e
z_S$ and $d_v z_S$, and the scalar recipes $m^2,V'_0(y)w,m^4$. Their
normalized polynomial envelopes have degree at most six. The sixth moment is
needed only to leading order, for which the first localization and a crude
perturbation bound suffice; no differentiated degree-eight recipe is required.

Removing the joint clips costs, by
\eqref{cap:joint-tail-gain},
\[
 |\langle\mathbf1_Ee z_S\rangle|
      \le Cs^{13/5}T_s^C,\qquad
 |\langle\mathbf1_Ed_v z_S\rangle|
      \le Cs^{21/10}T_s^C.
\]
We used $|e|\le C|y|$, boundedness of $m,V,A$, and $|\dot r|\le C\sqrt s$. The
sparse $A-V$ term obeys the second bound by the same joint row-energy product.
Scalar unclipping uses $m^2\le y^2$ and $|V'_0(y)|\le2|y|$; its products obey
the identical ledger. Gaussian unclipping has superpolynomially small tails
because the normalized surrogate variances are bounded.

We have now bounded the comparison, Taylor, and unclipping errors. The
Gaussian variables \(Y,G,U\) in the remaining calculation are centered and
have the current empirical Gram matrix of \(y,w,z_S\); they are the
full-frame surrogates, rather than the noncentral scalar equality laws used
earlier. The scalar comparisons and the justified replacement of \(A\) by
\(V\) give the precise target errors
\[
 |q-\E\tanh^2Y|=o(s^2),\qquad
 |\dot r-\E[V'_0(Y)G]|=o(s^{3/2}).
\]
Indeed the natural scalar sizes are \(s\) and \(\sqrt s\), and their
comparison and unclipping errors gain more than one power of \(s\).
Gaussian integration by parts now gives
\[
 \begin{split}
 \E[(\tanh Y-rY)U]
   &=(q-\E\tanh^2Y)\E[YU],\\
 \E[((V_0(Y)-r)G-\dot rY)U]
   &=(q-\E\tanh^2Y)\E[GU]
     +(\E[V'_0(Y)G]-\dot r)\E[YU].
 \end{split}
\]
The Gram bounds are \(|\E[GU]|\le C\) and \(|\E[YU]|\le C\sqrt s\).
The two displays therefore give \(o(s^{5/2})\) and \(o(s^2)\),
respectively. Every entry of
\eqref{cap:boundary-power-ledger} is strictly greater than one, and each
unclipping exponent is strictly above its target. By \eqref{cap:clip-choice}
the finite clipping losses leave positive margins. The supremum over unit
$S$-tests proves \eqref{cap:low-source-cancellation}. Quantitatively, after
the fixed-scale dimensional errors are decreased, its two right-hand sides are
bounded by $Cs^{5/2+1/200}$ and $Cs^{2+1/200}$, respectively; the smallest
remaining margin is $\zeta-1/1000>1/200$. This also proves the scalar
comparison for $m^4$ with error $o(s^3)$ and for $m^6$ to leading order.

\paragraph{Separation, stability, and response.}
Suppose \(q\le\varepsilon\sqrt t\), where \(\varepsilon>0\) will be fixed
below the field lower-bound constant.
The exact field equation is
\[
 h=r(z-W)y+(r-W)e.
\]
On the band \(2-W\le t\), its first multiplier is at most \(C(t+q^2)\)
in norm. Also \(s\asymp\sqrt t\) under the supposition, and
\(t\le s^2\le s^{7/5}=a\) at small scale, so
\(P_t^W\preceq P_S\). The norm bound for \(y\) and
\eqref{cap:low-source-cancellation} therefore give
\[
 |P_t^Wh|_n
 \le\bigl(C\sqrt\varepsilon+o_t(1)+o_n(1)\bigr)t^{5/4}.
\]
Choose \(\varepsilon\) so that \(C\sqrt\varepsilon\) is below half the
prescribed field lower-bound constant, then decrease the upper boundary
scale, and finally the dimensional error. This contradicts the field lower
bound, proving \(q\ge c\sqrt t\). In particular \(s\) is now comparable to
\(q\), and the identity
$z-2=q^2/(1-q)$ gives the positive separation needed to use the full
resolvent: $R>0$ and $\|R\|\le Cq^{-2}$. Inverting on all bands now gives
$|y-r^{-1}Rh|_n=o(\sqrt q)$ and $|m-Rh|_n=o(\sqrt q)$. Replacing $z$ by
$2+q^2$ incurs only a relative $o(1)$ error; consequently $q$ approximately
solves the scalar equation
\[
 q=|(2+q^2-W)^{-1}h|_n^2.
\]
Its right-hand side is decreasing in \(q>0\); equivalently,
\(q-|(2+q^2-W)^{-1}h|_n^2\) has derivative at least one.

For derivative tests, use \eqref{cap:derivative-band-equation},
\eqref{cap:first-derivative-localization}, and
\eqref{cap:low-source-cancellation} to obtain, uniformly under
\eqref{cap:derivative-equation},
\begin{equation}\label{cap:rank-response}
 w=r^{-1}R_\ell h_v
        -4qR_\ell y\langle yw\rangle+o(1)
 \quad\text{in normalized norm}.
\end{equation}
To see the coefficient, substitute $v=rw+\dot r\,y+d_v$ into
\eqref{cap:derivative-equation} and use $r-r^{-1}=-2q+O(q^2)$, $\dot
r=-2\langle yw\rangle+o(\sqrt q)$. The unenhanced term $\dot r\,y$ is $O(q)$,
whereas the resolvent term has coefficient $2q\dot r=-4q\langle
yw\rangle+o(q^{3/2})$. The low source errors are multiplied by at most
$Cq^{-2}$; their strict power gain therefore leaves an $o(1)$ error. The
high-band errors are controlled by the preceding localization estimates. This
justifies the remainder in the rank-response formula uniformly over the tests.

The rank-one operator $I+4qR_\ell yy^{\mathsf T}/n$ has uniformly bounded
inverse: its denominator is at least one, and its numerator correction has
norm at most $4q\|R_\ell\||y|_n^2\le C$. Thus \eqref{cap:rank-response} with
zero source excludes a normalized unit solution. If the stability matrix had a
nonpositive eigenvalue, increase $r-W$ by $\ell I$ until its minimum
eigenvalue reaches zero. The norm bound, $A\le(1+\eta)I$, and $q=O(s)$ ensure
$0\le\ell\le Cs$, choosing $\eta$ sufficiently small at the fixed bottom
scale. For invertible $A$, writing its eigenvector as $A^{1/2}w$ produces
exactly the zero-source test above, a contradiction. Likewise a positive
eigenvalue $\varepsilon\ll q^2$ at $\ell=0$ produces the source
$h_v=\varepsilon w$ and contradicts the same inverse estimate. Slightly
smaller tolerances and continuity include noninvertible \(A\). Hence the
stability test has positive fixed slack. The deterministic root lemma now
gives a unique root, and the implicit function theorem defines its response
on this stable field family. For \(A=V\) and arbitrary sources of
size $O(q^2)$, \eqref{cap:rank-response} first bounds the response and then,
by the rank-one inverse formula, gives \eqref{cap:tap-leading};
$|(I-V)w|_n=o(1)$ by the localization just proved.

\paragraph{The value of the potential.}
The root and response estimates determine the local curvature. To build a
single extension at initialization we also need to compare values at
separated fields. We therefore expand the stationary TAP potential through
cubic order in the overlap. The scalar fourth and sixth moments come from
the same Gaussian comparison, while the quadratic source correction uses
the low-band cancellation and the signed high-rotation estimate already
proved. Their coefficients must be combined before discarding the lower
powers.

The stationary TAP expression has gradient $m$. Since $h=(z-W)m-e/r$,
completing the square gives the first equality in \eqref{cap:tap-value}
exactly. Joint tail control permits the entropy expansion
\[
 \langle I(m)\rangle
 =q/2+\langle m^4\rangle/12+\langle m^6\rangle/30+o(q^3).
\]
The scalar comparisons just established express these moments by their
Gaussian Taylor expansions.  If $\tau=|y|_n^2$ denotes the matching
variance, then
\[
 \tau=q+2q^2+o(q^2),\quad
 \langle m^4\rangle=3q^2-8q^3+o(q^3),\quad
 \langle m^6\rangle=15q^3+o(q^3).
\]
For the quadratic correction, the immediate goal is
\(\langle eR'e\rangle=o(q^3)\). The low spectral region contributes at
most \(Cq^{-2}|P_Se|_n^2=o(q^3)\) by
\eqref{cap:low-source-cancellation}. On \(a<d\le u\), the clipped source
\(e_0\) contributes \(O(q^3u^{1/2}+q^5/a)T_s^C=o(q^3)\) by
\eqref{cap:clipped-root-source}.

Keep the coordinate-sparse source \(e_{\rm sp}=e-e_0\) separate on all
bands \(d>a\). Its improved row norm is
\(|e_{\rm sp}|_n\le Cq^{31/20}T_s^C\). The bounded sparse compression of
\(R_{\rm hi}\), together with \(R'=r^{-2}R-r^{-1}I\), bounds its
quadratic form by \(O(q^{31/10}T_s^C)=o(q^3)\). The bandwise sparse estimate
following \eqref{cap:first-root-localization} and the improved row bound
make its cross terms on the middle bands \(o(q^3)\). Here the spectral projection
belongs to the operator: for example,
\(\langle P_{>u}e_{\rm sp},R'P_{>u}e_{\rm sp}\rangle
=\langle e_{\rm sp},R'_{>u}e_{\rm sp}\rangle\). The vector tested by the
sparse estimate is still \(e_{\rm sp}\).

Above \(u\), write \(e_0=t_y+e_{\rm rep}\), where
\(e_{\rm rep}=e_0(y)-e_0(Y_0)\); the definition of \(t_y\) in
\eqref{cap:low-input-recipes} uses the same cutoff \(\chi\).
For the structured term,
\(|t_y|_n\le Cq^{3/2}T_s^C\), and the signed high test gives
\[
 |\langle t_yR'_{>u}t_y\rangle|
 \le Cq^{1/20}T_s^C|t_y|_n^2=o(q^3).
\]
The high-input replacement estimates give
\(|R'_{>u}e_{\rm rep}|_n\le Cq^{15/8}T_s^C\) and
\(|e_{\rm rep}|_n\le Cq^{79/40}T_s^C\). Its cross term with \(t_y\)
is therefore \(O(q^{27/8}T_s^C)\), and its self term is of still higher
order. For the cross term with \(e_{\rm sp}\), the structured sparse-output
bound gives
\(|\mathbf1_E R'_{>u}t_y|_n\le Cq^{31/20}T_s^C\), while the full
replacement bound above handles \(R'_{>u}e_{\rm rep}\). Multiplication by
\(|e_{\rm sp}|_n\) gives powers \(31/10\) and \(137/40\), respectively.
All exceed three after the clipping losses. This proves
\(\langle eR'e\rangle=o(q^3)\). It follows from \(R=r^2R'+rI\) that
\[
 \langle eRe\rangle=r\langle e^2\rangle+o(q^3)
                         =\tfrac23q^3+o(q^3).
\]
The last scalar coefficient follows by expanding $e=qY-Y^3/3+O(Y^5)$ under the
matched Gaussian law. Substitution cancels the quadratic terms and proves
\eqref{cap:tap-value}, with $\kappa_0=1/3$. This coefficient is used in
the own-saddle support comparison at initialization.

Residuals at approximate roots and changes of coefficient or matrix add small
normalized sources to the two exact equations used above. At a fixed bottom
scale all inverses, clip derivatives, and tests have fixed bounds. Choose
these perturbations after the displayed scales and error targets, and repeat
the estimates with their reserved slack. The TAP expression then uses the
actual $J_*$ and $jr^2/4$. This proves the final perturbation assertion and
completes the proof. \end{proof}
